\documentclass{article}

\usepackage[T2A]{fontenc}
\usepackage[utf8]{inputenc}
\usepackage[russian]{babel}

\usepackage[unicode, colorlinks, linkcolor=blue, citecolor=blue]{hyperref}
\usepackage{amsmath}
\usepackage{graphicx}
\graphicspath{{pictures/}}
\DeclareGraphicsExtensions{.pdf,.png,.jpg}

\usepackage{pgfplots}
\usepackage{tikz}


\begin{document}
	\selectlanguage{russian}
	%----------------------------------Шапка-------------------------------------------
	\begin{center}
		\includegraphics[scale=0.25]{AU}\\
		{\Large\bfseries Санкт-Петербургский национальный исследовательский Академический университет имени Ж.И.~Алфёрова Российской академии наук}
	\end{center}
	
	\begin{center}
		{\large\textbf{Рабочий протокол и отчёт по лабораторной работе № 5}}\\
		Свиридов Фёдор, Александр Слободнюк, Владимир Попов
	\end{center}
	
	\begin{center}
		\rule{12cm}{0.4mm}\\
		\large\bfseries{<<Проверка закона Шарля>>}\\
		\rule{12cm}{0.4mm}
	\end{center}
	%---------------------------------------------------------------------------------------
	\paragraph{Цель работы.} Исследование изохорного процесса
	\paragraph{Задачи, решаемые при выполении работы.}
	\begin{enumerate}
		\item Нагреть газ до 50 градусов цельсия
		\item Начать охлаждение газа
		\item Уменьшать объем, занимаемый газом, после понижения давления на определенную величину
		\item Измерять температуру и давление при этой температуре
		\item Сделать выводы
	\end{enumerate}
	\paragraph{Объект исследования.} Модель идеального газа
	\paragraph{Метод экспериментального исследования.} Поэтапное измерение температуры
	\paragraph{Исходные данные.}
	Будем считать воздух идеальным газом, тогда:
	$$PV=\nu RT $$
	$$ P=\frac{\nu R}{V}\,T$$
	В нашем опыте количество вещества $\nu$ оставалось примерно постоянным, а вот объём $V$ нам приходилось немного изменять. Опишем несколько последовательных состояний нашей системы.
	Пусть в самом начале опыта система находилась в состоянии $F(P_0,V_0,T_0)$
	$$ P_0=\frac{\nu R}{V_0}\,T_0\; \xrightarrow{(1)} \; P_1=\frac{\nu R}{V_0}\,T_1 \;\xrightarrow{(2)}\; P_0=\frac{\nu R}{V_0+dV}\,T_1 \;\xrightarrow{(3)} \; P_1=\frac{\nu R}{V_0+dV}\,T_2$$
	
	(1) - изохорный процесс с коэффициентом $ \frac{\nu R}{V_0} $
	
	(2) - возврат к давлению $P_0$ с помощью изменения объёма
	
	(3) - изохорный процесс с коэффициентом $ \frac{\nu R}{V_0+dV} $\\
	
	Таким образом, $\Delta P=C(V)\Delta T$, где C(V) - некоторый коэффициент пропорциональности, который зависит от объёма. Но если пренебречь величиной $dV$, то можно считать, что $\Delta P \sim \Delta T$\\
	
	\textbf{Результаты прямых измерений и их обработки.}
	
	\begin{figure}[htb]
		\begin{minipage}{0.5\linewidth}
			\centering{Нагревание воздуха}\\
			\begin{tabular}{c|c}
				$\Delta P$, кПа&$\Delta T$, $^\circ\mbox{C}$\\
				\hline
				0,2&5,6 \\
				0,2&3,7  \\
				0,2&2,4 \\
				0,2&2,7 \\
				0,2&3,0 \\
				0,2&2,9 \\
			\end{tabular}
		\end{minipage}
		\hfill
		\begin{minipage}{0.5\linewidth}
			\centering{Остывание воздуха}\\
			\begin{tabular}{c|c}
				$\Delta P$, кПа&$\Delta T$, $^\circ\mbox{C}$\\
				\hline
				0,2&2,0 \\
				0,2&1,7  \\
				0,2&1,5  \\
				0,2&1,0  \\
				0,2&1,0 \\
				0,2&1,1 \\
				0,2& 1,2 \\
				0,2& 1,0 \\
			\end{tabular}
		\end{minipage}
	\end{figure}
	
	Найдём зависимость $ P(t) $ на основе данных остывания воздуха. В начальный момент $ t_0=49,6\; ^{\circ}C $, $ P_0=756 \;\mbox{мм. рт. ст.} $ или $ P_0=100{,}775\; $кПа
	
	\begin{figure}[htb]
		\centering\caption{Зависимость давления от температуры}
		\begin{tikzpicture}[scale=0.9]
			\begin{axis}[legend pos = north west,xlabel={Температура $t$, $ ^{\circ}C $}, ylabel={Давление $ P{,}\;$кПа}]
				\legend{,$ 0.151311\,t+93.3654 $};
				
				\addplot [ only marks, mark=*, mark size=1.5pt,  error bars/.cd, y dir=both, y explicit]
				table [x=x, y=y, y error=y-err]
				{
					x     y  y-err
					49.6 	100.775  	0.2
					47.6	 100.575 	0.2
					45.9	100.375		0.2
					44.4	 100.175	0.2
					43.4	99.975		0.2
					42.4	 99.775		0.2
					41.3	99.575		0.2
					40.1	99.375		0.2
					39.1	99.275		0.2
				};
				
				\addplot [draw=red][domain=39:50]{93.3654 + 0.151311*x};
			\end{axis}
		\end{tikzpicture}
	\end{figure}
	
	\begin{figure}[htb]
		\centering{Координаты точек}\\
		
		\begin{tabular}{c|c}
			$t$, $^\circ\mbox{C}$&$P$, кПа\\
			\hline
			49.6&100.775\\
			47.6& 100.575\\
			45.9&100.375\\
			44.4& 100.175\\
			43.4& 99.975\\
			42.4& 99.775\\
			41.3&99.575\\
			40.1&99.375\\
			39.1&99.275\\
		\end{tabular}
	\end{figure}
	
	Пусть $P(0)=93.4$ кПа - давление газа при температуре $0\;^\circ\mbox{C}$. Тогда экстраполяционная прямая принимает вид:
	$P(0)(1+0,0016\,t)$. \fbox{$k = 0,0016$}\\
	Решение уравнения $  0.151311\,t+93.3654=0$ даёт следующий результат: $ t=-617\;^\circ\mbox{C} $\\
	
	
	\textbf{Выводы и анализ результатов.}
	Мы провели измерения изменения давления $\Delta P$ и температуры $\Delta T$, чтобы проверить закон Шарля. На основе полученных данных была получена линейная зависимость~$P \sim t$.\\
	Полученный коэффициент наклона прямой $k = 0,0016$ по порядку совпадает с ожидаемой величиной $\frac{1}{273}$. Такой большой промах $ t=-617\;^\circ\mbox{C} $ температуры абсолютного нуля можно объяснить огромной погрешностью в данном опыте, которую мы не стали оценивать. Полученные результаты не очень убедительны, потому что опыт обладает рядом недостатков. Существенной проблемой является то, что установка при повышенном давление пропускает воздух, таким образом, невозможно качественно провести проверку закона Шарля: при постоянном объёме и количестве вещества нагревать (охлаждать) газ и следить за повышением (понижением) давления. Следующим недостатком опыта является малая точность манометра, из-за чего возникают трудности при выравнивание давления до начального состояния $P_0$. Из несущественных недостатков можно отметить, что воздух является не идеальным газом и то, что нам приходилось менять объём на небольшую величину, которой мы пренебрегли.
	
\end{document}